% =====================================================================
% Seção 5 — Discussão
% =====================================================================
\section{Discussão}\label{sec:disc}

A convergência entre a literatura e o estudo de caso permite articular o fator
limítrofe em três planos complementares, cada um ancorado em evidência.

\subsection{Plano informacional: a informação não se auto-cria}

A evidência mais forte é informacional. O colapso sob treinamento recursivo
(\emph{Nature}, 2024)~\parencite{Shumailov2024Collapse} e sua inevitabilidade
estatística~\parencite{Borji2024CollapseNote} mostram que \emph{informação nova
não é criada por reamostragem}. A derivação teórica do expoente de scaling a
partir das estatísticas da fonte~\parencite{Cagnetta2026ScalingTheory} reforça
esse ponto: o ``teto'' do loop é a entropia da fonte externa. No estudo de caso,
o análogo é direto: o HABD reorganiza âncoras já existentes e obtém diversidade,
mas não expande o conteúdo informacional do corpus; quando pressiona a métrica,
o custo aparece em groundedness.

\subsection{Plano de avaliação: gerador e avaliador não podem coincidir}

O segundo plano é o sinal de avaliação. A literatura de
\emph{self-play}~\parencite{SelfPlayGuidance2026} e o
\emph{discovery--reliability gap}~\parencite{RSIBench2026} mostram que, quando o
mesmo processo gera e avalia, o feedback perde validade. No estudo de caso, a
\texttt{groundedness} atua como avaliador \emph{independente} do gerador de
diversidade --- é ela que impede o overclaim. Este é o comportamento que a
definição operacional da Seção~\ref{sec:metodo} exige (condição (ii)): sem um
sinal externo de fidelidade, o ganho de diversidade seria reportado como vitória
limpa.

\subsection{Plano de mitigação: âncora real, não teto absoluto}

O terceiro plano distingue o teto de um loop \emph{fechado} do de um loop
\emph{ancorado}. Gerstgrasser \textit{et al.}~\parencite{Gerstgrasser2024Collapse}
e Dohmatob \textit{et al.}~\parencite{Dohmatob2024CollapseScaling} mostram que o
colapso é \emph{contido} quando dados reais são acumulados a cada iteração. Isso
se transfere ao estudo de caso e, por extensão, ao \emph{OpenCode Ecosystem Core}
em geral: cada ciclo de pesquisa (R458--R460) é um loop ancorado --- a
diversificação de código, a geração de dados sintéticos para testes e a
supervisão por agentes auditáveis funcionam como âncoras que impedem o loop de
degenerar. O ecossistema, ao exigir especificação formal e testes (gate SDD/TDD)
antes de qualquer entrega, incorpora praticamente a condição (ii) da definição:
o avaliador (teste) é externo ao gerador (código).

\subsection{Conexão com o contexto do OpenCode Ecosystem Core}

O estudo de caso não é o único lugar do ecossistema onde o fator limítrofe emerge.
O \emph{OpenCode Ecosystem Core} opera 278 ciclos de evolução registrados, 124
especificações formais e um MetaBus de memória metacognitiva compartilhada. Nesse
desenho, os agentes geram artefatos (código, textos, análises), mas a \emph{validação
} é sempre delegada a uma instância distinta (gate SDD/TDD, BehavioralGate,
auditoria bibliográfica, revisão cega emulada). A \emph{razão} dessa arquitetura é
precisamente a hipótese que este artigo defende: \textbf{produzir IA com IA é
viável e produtivo quando há âncora externa de validação; torna-se degenerativo
quando o mesmo processo gera e julga}. O ecossistema é, nesse sentido, um
experimento permanente do fator limítrofe --- e seus vereditos honestos
(\texttt{refuta\_H2}, \texttt{refuta\_H3}) são exemplos de \emph{anti-overclaim}
operacionalizado.

\subsection{Limites desta evidência}

É preciso honestidade epistêmica. Grande parte das fontes de auto-melhoria são
preprints (arXiv) ainda não revisados por pares; apenas o núcleo --- colapso
(\emph{Nature} 2024), scaling laws (Kaplan 2020, Chinchilla 2022) --- é
peer-reviewed consolidado. O estudo de caso usa \textbf{4 consultas} em
\emph{corpus-piloto controlado}; seu papel é \emph{ilustrativo e internalmente
válido}, não estatisticamente generalizável. Não afirmamos que o fator limítrofe
seja inviolável para sempre: melhorias algorítmicas podem elevar o teto, e a
transferência dos resultados de distribuições simuladas para linguagem real é
objeto de debate legítimo (cf.~emergência como artefato de métrica,~\textcite
{Schaeffer2023EmergentMirage}).
